\subsection{Smooth supports and four discriminants}
\label{ssec:smoothsupport}

\Cref{thm:candidatefails} and \Cref{cor:localobstruction} exclude every secant
object whose support has two transversal branches of codimension two
through a point. This leaves supports that are smooth or that meet
themselves only along curves, and \Cref{rem:candidatemeaning} ends with the
observation that no such object is known. In this subsection we explain why
the local argument cannot reach the smooth case, and then determine what a
smooth support can be. The answer is narrow: on the fourfold the invariants
of the support are forced, and the discriminant is confined to a finite list.

We begin with the fact that at a smooth point there is nothing local left to
compute.

\begin{proposition}[No local obstruction at a smooth point]\label{prop:smoothlocal}
Let $Y$ be a smooth variety with $\omega_{Y}\cong\cO_{Y}$ and let $S\subset Y$
be smooth of codimension two. Then
\[
  \cE xt^{0}(\cI_{S},\cI_{S})=\cO_{Y},\qquad
  \cE xt^{1}(\cI_{S},\cI_{S})=N_{S/Y},\qquad
  \cE xt^{q}(\cI_{S},\cI_{S})=0\quad(q\ge2),
\]
and the same holds for $\cI_{S}\otimes L$ for any line bundle $L$.
Consequently, for $F=\cI_{S}\otimes L$ both sides of \eqref{eq:heisenberg}
have vanishing germs in $\cE xt^{2}$, and \Cref{lem:edge} decides nothing.
\end{proposition}

\begin{proof}
Twisting by $L$ changes neither side, so take $F=\cI_{S}$. The sheaf
$\cI_{S}$ is torsion free of rank one, and its dual is $\cO_{Y}$ because $S$
has codimension two; hence every local endomorphism of $\cI_{S}$ is
multiplication by a function, and
$\cE xt^{0}(\cI_{S},\cI_{S})=\cO_{Y}$. Apply $\cH om(-,\cI_{S})$ to
\begin{equation}\label{eq:idealseq}
  0\to\cI_{S}\to\cO_{Y}\to\cO_{S}\to0 .
\end{equation}
Since $\cE xt^{q}(\cO_{Y},\cI_{S})=0$ for $q\ge1$, the long exact sequence
gives
\begin{equation}\label{eq:shift}
  \cE xt^{q}(\cI_{S},\cI_{S})\cong\cE xt^{q+1}(\cO_{S},\cI_{S})
  \qquad(q\ge1).
\end{equation}
Now apply $\cH om(\cO_{S},-)$ to \eqref{eq:idealseq}. Because $S$ is smooth of
codimension two, hence a local complete intersection, the Koszul resolution
gives $\cE xt^{k}(\cO_{S},\cO_{Y})=0$ for $k\ne2$ and
$\cE xt^{2}(\cO_{S},\cO_{Y})=\omega_{S}\otimes\omega_{Y}^{-1}=\omega_{S}$,
while $\cE xt^{k}(\cO_{S},\cO_{S})=\bigwedge^{k}N_{S/Y}$ for $0\le k\le2$ and
zero otherwise. Computing both from the same Koszul complex shows that the
map
\[
  \rho\colon\cE xt^{2}(\cO_{S},\cO_{Y})=\omega_{S}
  \longrightarrow
  \cE xt^{2}(\cO_{S},\cO_{S})=\textstyle\bigwedge^{2}N_{S/Y}=\omega_{S}
\]
induced by $\cO_{Y}\to\cO_{S}$ is the identity: in local coordinates in which
$S=\{x_{1}=x_{2}=0\}$, both groups are computed as $M/(x_{1},x_{2})M$ for
$M=\cO_{Y}$ and $M=\cO_{S}$, and the map is reduction modulo $(x_{1},x_{2})$,
which is the identity on $\cO_{S}$.

Consider the long exact sequence of $\cH om(\cO_{S},-)$ in degrees two and
three. Since $\rho$ is surjective and $\cE xt^{3}(\cO_{S},\cO_{Y})=0$, the
segment
\[
  \cE xt^{2}(\cO_{S},\cO_{Y})\xrightarrow{\ \rho\ }
  \cE xt^{2}(\cO_{S},\cO_{S})\to
  \cE xt^{3}(\cO_{S},\cI_{S})\to
  \cE xt^{3}(\cO_{S},\cO_{Y})=0
\]
gives $\cE xt^{3}(\cO_{S},\cI_{S})=0$, so $\cE xt^{2}(\cI_{S},\cI_{S})=0$ by
\eqref{eq:shift}. For $k\ge3$ both $\cE xt^{k}(\cO_{S},\cO_{Y})$ and
$\cE xt^{k}(\cO_{S},\cO_{S})$ vanish, so
$\cE xt^{k+1}(\cO_{S},\cI_{S})=0$ and $\cE xt^{k}(\cI_{S},\cI_{S})=0$. In
degrees one and two, since $\rho$ is injective, the segment
\[
  0=\cE xt^{1}(\cO_{S},\cO_{Y})\to
  N_{S/Y}\to
  \cE xt^{2}(\cO_{S},\cI_{S})\to
  \cE xt^{2}(\cO_{S},\cO_{Y})\xrightarrow{\ \rho\ }
  \cE xt^{2}(\cO_{S},\cO_{S}),
\]
in which $N_{S/Y}=\cE xt^{1}(\cO_{S},\cO_{S})$, gives
$\cE xt^{2}(\cO_{S},\cI_{S})\cong N_{S/Y}$, and this is
$\cE xt^{1}(\cI_{S},\cI_{S})$ by \eqref{eq:shift}.

For the last assertion, the right side of \eqref{eq:heisenberg} has vanishing
germs by \Cref{lem:edge}, and the left side lies in
$\Ext^{2}(F,F\otimes\Omega^{2}_{Y})$, whose sheaf of germs in the
$\cE xt^{2}$ direction is $\cE xt^{2}(F,F)\otimes\Omega^{2}_{Y}=0$.
\end{proof}

\begin{remark}[The global form of the criterion]\label{rem:obstructionmoves}
\Cref{prop:smoothlocal} is the precise reason why the argument of
\Cref{thm:candidatefails} stops at a smooth support, and it also shows what
takes its place. Since $\cE xt^{2}=0$, the local-to-global spectral sequence
for $\Ext^{2}(\cI_{S},\cI_{S})$ has only two surviving rows,
$H^{1}(S,N_{S/Y})$ and $H^{2}(Y,\cO_{Y})$, and they sit in different steps of
the filtration: the first is the quotient that records the geometry of $S$,
the second the subgroup of scalars. A product $A_{\theta}A_{\theta'}$ of
translation classes therefore has a well defined image in $H^{1}(S,N_{S/Y})$,
while the scalar on the right of \eqref{eq:heisenberg} has image zero there.
The criterion thus splits into two conditions of different kinds: a vanishing
in $H^{1}(S,N_{S/Y})$, which is a statement about the normal bundle and not
about any point, and an identity in $H^{2}(Y,\cO_{Y})$. At a double point the
first condition already fails locally. At a smooth support it is a global
question, and what is decided here is which surfaces $S$ are available at
all.
\end{remark}

The invariants of a smooth support are far from free. Since the Todd class of
an abelian variety is trivial, Grothendieck--Riemann--Roch converts the Chern
character of the object into the complete set of numerical invariants of the
surface, and three classical inequalities then constrain them.

\begin{theorem}[The invariants of a smooth support]\label{thm:smoothinvariants}
Let $(X,\Theta)$ be a principally polarised abelian fourfold, let $d>0$, let
$b\in\ZZ_{>0}$, and let $\cF=\cI_{S}(b\Theta)$ be a coherent sheaf with
\[
  \ch(\cF)=u_{\Theta}+b\,v_{\Theta}
\]
and $S\subset X$ smooth. Write $N=\tfrac12\Nm_{K/\QQ}(b+\sqrt{-d})
=\tfrac12(b^{2}+d)$, $\iota\colon S\hookrightarrow X$ and
$\lambda=\iota^{*}\Theta$. Then $S$ is a surface and
\begin{equation}\label{eq:smoothinvariants}
\begin{aligned}
  [S]&=N\Theta^{2}, &\qquad \iota_{*}K_{S}&=\tfrac{4bN}{3}\,\Theta^{3},
  &\qquad \chi(\cO_{S})&=4N(2b^{2}-N),\\
  K_{S}^{2}&=12N(4b^{2}-N), &\qquad e(S)&=12N(4b^{2}-3N),
  &\qquad K_{S}\cdot\lambda&=32bN,
\end{aligned}
\end{equation}
with $\lambda^{2}=24N$ and $[S]^{2}=24N^{2}$. Moreover $K_{S}$ is globally
generated, hence nef; $S$ is minimal of general type; and
\begin{equation}\label{eq:threebounds}
  \tfrac{4}{9}b^{2}\ \le\ N\ \le\ b^{2},
  \qquad\text{equivalently}\qquad
  -\tfrac{1}{9}b^{2}\ \le\ d\ \le\ b^{2} .
\end{equation}
The lower bound is the Hodge index theorem and is implied by $d>0$; the upper
bound is the Bogomolov--Miyaoka--Yau inequality. In terms of $b$ and $d$ the
five invariants read
\begin{equation}\label{eq:bdform}
  \chi(\cO_{S})=(b^{2}+d)(3b^{2}-d),\quad
  K_{S}^{2}=3(b^{2}+d)(7b^{2}-d),\quad
  e(S)=3(b^{2}+d)(5b^{2}-3d),
\end{equation}
and the defect in the Bogomolov--Miyaoka--Yau inequality is
\begin{equation}\label{eq:bmydefect}
  3e(S)-K_{S}^{2}=24\,(b^{4}-d^{2}),
\end{equation}
so that the inequality amounts exactly to $d\le b^{2}$.
Finally, $K_{S}$ is never numerically proportional to $\lambda$.
\end{theorem}

\begin{proof}
Since $\ch_{0}(\cF)=1$ and $c_{1}(\cF)=b\Theta$, the sheaf
$\cF(-b\Theta)=\cI_{S}$ has $\ch_{0}=1$, $\ch_{1}=0$, so $S$ has codimension
at least two, and
\[
  \ch(\cO_{S})=1-\ch(\cI_{S})=1-\ch(\cF)\,e^{-b\Theta}.
\]
Put $z=b+\sqrt{-d}$, so that $z\bbar z=2N$. Since
$u_{\Theta}=\Re e^{\sqrt{-d}\Theta}$ and
$\sqrt d\,v_{\Theta}=\Im e^{\sqrt{-d}\Theta}$, we have
$\ch(\cF)=\Im\bigl(z\,e^{\sqrt{-d}\Theta}\bigr)/\sqrt d$, hence
$\ch(\cF)e^{-b\Theta}=\Im\bigl(z\,e^{-\bbar z\Theta}\bigr)/\sqrt d$, whose
coefficient of $\Theta^{k}$ is
$(-1)^{k}\Im(z\bbar z^{k})/(\sqrt d\,k!)=(-1)^{k+1}\,2N\Im(z^{k-1})/(\sqrt d\,k!)$.
With $\Im z=\sqrt d$, $\Im z^{2}=2b\sqrt d$ and
$\Im z^{3}=(3b^{2}-d)\sqrt d$ this gives
\[
  \ch_{2}(\cO_{S})=\tfrac{b^{2}+d}{2}\Theta^{2}=N\Theta^{2},\quad
  \ch_{3}(\cO_{S})=-\tfrac{2bN}{3}\Theta^{3},\quad
  \ch_{4}(\cO_{S})=\tfrac{3b^{4}+2db^{2}-d^{2}}{24}\Theta^{4}.
\]
The first
identity says $S$ has codimension exactly two, so it is a surface, with class
$N\Theta^{2}$. Because $T_{X}$ is trivial, $\td(T_{X})=1$, and
Grothendieck--Riemann--Roch for $\iota$ reads
$\ch(\iota_{*}\cO_{S})=\iota_{*}\td(T_{S})$. Comparing degrees three and four,
\[
  \ch_{3}(\cO_{S})=\iota_{*}\tfrac{c_{1}(T_{S})}{2}=-\tfrac12\iota_{*}K_{S},
  \qquad
  \ch_{4}(\cO_{S})=\iota_{*}\tfrac{c_{1}^{2}+c_{2}}{12}
  =\chi(\cO_{S})\,[\mathrm{pt}],
\]
the second by Noether's formula. With $\Theta^{4}=24\,[\mathrm{pt}]$ these give
$\iota_{*}K_{S}=\tfrac{4bN}{3}\Theta^{3}$ and, after substituting
$d=2N-b^{2}$, $\chi(\cO_{S})=3b^{4}+2db^{2}-d^{2}=4N(2b^{2}-N)$.
Then $K_{S}\cdot\lambda=\int_{X}\iota_{*}K_{S}\cdot\Theta=\tfrac{4bN}{3}\cdot24=32bN$
and $\lambda^{2}=\int_{X}[S]\cdot\Theta^{2}=24N$.

For the last two invariants, the normal bundle sequence and the triviality of
$T_{X}$ give $c(T_{S})c(N_{S/X})=1$, so $c_{1}(N_{S/X})=K_{S}$ and
$c_{2}(N_{S/X})=K_{S}^{2}-e(S)$; the self-intersection formula
$[S]^{2}=\int_{S}c_{2}(N_{S/X})$ then reads
\[
  24N^{2}=K_{S}^{2}-e(S),
\]
while Noether's formula reads $12\chi(\cO_{S})=K_{S}^{2}+e(S)$. Solving the
two gives the stated $K_{S}^{2}$ and $e(S)$.

The conormal sequence presents $\Omega^{1}_{S}$ as a quotient of
$\Omega^{1}_{X}|_{S}=\cO_{S}^{\oplus4}$, so $\Omega^{1}_{S}$ is globally
generated, and therefore so is $K_{S}=\det\Omega^{1}_{S}$, a quotient of
$\bigwedge^{2}\cO_{S}^{\oplus4}$. A globally generated line bundle is nef. A
globally generated rank two bundle on a surface has $c_{2}\ge0$, the class
$c_{2}$ being represented by the vanishing scheme of a general section, which
is empty or finite; hence
\[
  e(S)=c_{2}(\Omega^{1}_{S})=12N(4b^{2}-3N)\ \ge\ 0,
\]
so $N\le\tfrac43b^{2}<4b^{2}$ and therefore $K_{S}^{2}=12N(4b^{2}-N)>0$. A nef
divisor of positive self-intersection is big, so $\kappa(S)=2$ and, $K_{S}$
being nef, $S$ is minimal of general type. The Bogomolov--Miyaoka--Yau
inequality $c_{1}^{2}\le3c_{2}$ therefore applies and reads
$K_{S}^{2}\le3e(S)$. Substituting $N=\tfrac12(b^{2}+d)$ in
\eqref{eq:smoothinvariants} gives \eqref{eq:bdform}, hence
$3e(S)-K_{S}^{2}=24(b^{4}-d^{2})$, and the inequality is $d\le b^{2}$, that
is $N\le b^{2}$.

For the lower bound, $\lambda$ is ample on $S$, so the Hodge index theorem
gives $(K_{S}\cdot\lambda)^{2}\ge K_{S}^{2}\,\lambda^{2}$, that is
$(32bN)^{2}\ge12N(4b^{2}-N)\cdot24N$, which simplifies to $9N\ge4b^{2}$. Since
$N=\tfrac12(b^{2}+d)$ with $d>0$, this holds automatically. Equality in the
Hodge index theorem occurs exactly when $K_{S}$ is numerically proportional to
$\lambda$, hence exactly when $9N=4b^{2}$, and this never happens, because
$9N=\tfrac92(b^{2}+d)>4b^{2}$ for $d>0$.
\end{proof}

\Cref{fig:bmywindow} draws the Bogomolov--Miyaoka--Yau defect over the
quarter plane of $(b,d)$.

\begin{corollary}[Four discriminants on the fourfold]\label{cor:fourdiscriminants}
Let $(X,\Theta)$ be a principally polarised abelian fourfold and $d>0$
squarefree. A sheaf $\cI_{S}(3\Theta)$ with Chern character
$u_{\Theta}+3v_{\Theta}$, the value of \Cref{prop:markman8}(iii), and with
smooth support $S$ exists only if $d\le9$, that is
$d\in\{1,2,3,5,6,7\}$; if in addition $N=(9+d)/2$ is an integer, as in
\Cref{setup:markman}, only
\[
  d\in\{1,\,3,\,5,\,7\},\qquad N\in\{5,\,6,\,7,\,8\}.
\]
For these values the invariants are forced to be those of
\Cref{tab:smoothsupport}. In particular, for $d>9$ no such sheaf has smooth
support.
\end{corollary}

\begin{proof}
This is immediate from \eqref{eq:threebounds} at $b=3$ and the arithmetic of
\eqref{eq:smoothinvariants}, and the table evaluates
\eqref{eq:smoothinvariants} at $b=3$. The value $d=9$ is the equality case of the
Bogomolov--Miyaoka--Yau inequality, where $K_{S}^{2}=3e(S)=2916$ and $S$ would
have to be a quotient of the ball; it is excluded because $9$ is not
squarefree.
\end{proof}

\begin{table}[htbp]
\centering
\small
\renewcommand{\arraystretch}{1.12}
\begin{tabular}{@{}ccrrrrr@{}}
\toprule
$d$ & $N$ & $\chi(\cO_{S})$ & $K_{S}^{2}$ & $e(S)$ &
$K_{S}\cdot\lambda$ & $\lambda^{2}$ \\
\midrule
$1$ & $5$ & $260$ & $1860$ & $1260$ & $480$ & $120$ \\
$3$ & $6$ & $288$ & $2160$ & $1296$ & $576$ & $144$ \\
$5$ & $7$ & $308$ & $2436$ & $1260$ & $672$ & $168$ \\
$7$ & $8$ & $320$ & $2688$ & $1152$ & $768$ & $192$ \\
\bottomrule
\end{tabular}
\caption{The only invariants a smooth support can have on an abelian
fourfold, by \Cref{thm:smoothinvariants} and
\Cref{cor:fourdiscriminants}, at $b=3$ and with $N$ an integer. Every entry is
determined by $\ch(\cF)$ alone. The slope
$K_{S}^{2}/\chi(\cO_{S})$ runs from $93/13$ at $d=1$ to $42/5$ at $d=7$, and
reaches the Bogomolov--Miyaoka--Yau ceiling of $9$ only at the excluded value
$d=9$.}
\label{tab:smoothsupport}
\end{table}

\begin{figure}[htbp]
\centering
  \includegraphics[width=\textwidth]{fig_bmywindow}
\caption{\Cref{thm:smoothinvariants} over the quarter plane of $(b,d)$. The
height of the surface is the defect $3e(S)-K_{S}^{2}=24(b^{4}-d^{2})$ in the
Bogomolov--Miyaoka--Yau inequality for the surface the theorem attaches to the
pair, compressed by a signed cube root so that four orders of magnitude fit
one frame; the compression is odd and monotone, so the zero locus and the sign
are the true ones. The surface is drawn indigo where the defect is positive and
clay where it is negative, and it crosses zero along $d=b^{2}$, which is
therefore the whole content of the inequality. The weaker wall $d=\tfrac53
b^{2}$, where $e(S)$ itself changes sign, is drawn dashed and lies above it.
On the line $b=3$, which is the value \Cref{prop:markman8}(iii) forces on the
fourfold, the four squarefree discriminants $d=1,3,5,7$ with $N$ an integer sit
above the wall and are marked with their drop lines; $d=9$ sits exactly on it,
the equality case, and $d=11,13,15$ sit below and are excluded.}
\label{fig:bmywindow}
\end{figure}

\begin{remark}[Status of the smooth case]\label{rem:smoothremains}
The local argument does not apply here, for a structural reason given by
\Cref{prop:smoothlocal}: the sheaf $\cE xt^{2}$ of a smooth ideal vanishes,
so both sides of \eqref{eq:heisenberg} are locally zero, and the comparison of
germs that decides \Cref{thm:candidatefails} has nothing to compare.

The surface is rigidly constrained. At $b=3$ its class, its canonical class,
its Euler characteristic, its Euler number and its two intersection numbers
with the polarisation are all determined by $d$, and they are compatible with
existence only for $d\le9$. This is the content of
\Cref{cor:fourdiscriminants}, and it restricts the construction: for all but
finitely many imaginary quadratic fields the secant construction of
\Cref{setup:markman} with smooth support is not available, whatever the
geometry of the surface.

What remains is a single existence question. For $d\in\{1,3,5,7\}$, does a
principally polarised abelian fourfold carry a smooth surface $S$ with the
invariants of \Cref{tab:smoothsupport}, and if so, does the image of
$A_{\theta}A_{\theta'}$ in $H^{1}(S,N_{S/X})$ vanish for all
$\theta,\theta'\in H^{0}(T_{X})$? The first half is a question about surfaces
in abelian fourfolds with prescribed Chern numbers, of the kind that is
answered by a construction or by a classification argument; the second is a
question about one cohomology group of one normal bundle. Neither is settled
here. The simplest surfaces of the right class are ruled out by one entry of
the table: a complete intersection of divisors in $|a\Theta|$ and
$|b'\Theta|$ has class $ab'\Theta^{2}$, so $ab'=N$, and canonical class
$(a+b')\lambda$ by adjunction, so $K_{S}\cdot\lambda=24(a+b')N$; since
$ab'=N\ge5$ forces $a+b'\ge5$, this exceeds the required $96N$. The
surface, if it exists, is not cut out by two theta divisors, nor by any two
divisors algebraically equivalent to multiples of $\Theta$; nor, by
\Cref{cor:latticeburch} below, is its ideal sheaf presented by any two bundles
whose classes lie in the subring generated by line bundles. It must also move
with $X$ to first order in every polarised direction, in the exact sense of
\Cref{prop:smoothrank}, while the injectivity that \Cref{lem:annihilation}
asks of it holds automatically when $X$ is simple (\Cref{lem:smoothinjective}).

A positive answer to both, for a surface that moves in this way, together
with an affirmative answer to \cite[Question 11.4]{Mar25b}, would produce
Weil classes on abelian
eightfolds for at most four discriminants; the other discriminants, every
$n\ge5$ and every field of degree larger than two would still have to be
reached through \Cref{prop:descent} and the propagation statement.
\end{remark}

Two of the conditions on a smooth support are decided by the table, and the
third, which \Cref{thm:factor}(iii)(a) imposes on every candidate, has an
exact form. The first is the injectivity of $H^{2}(\cO_{X})\to H^{2}(\cO_{S})$
that \Cref{lem:annihilation} below requires of every support.

\begin{lemma}[The injectivity condition for a smooth support]
\label{lem:smoothinjective}
Let $X$ be a simple complex abelian fourfold, let $V=H^{0}(T_{X})$, and let
$S\subset X$ be a smooth surface. The following are equivalent:
\begin{enumerate}[label=\textup{(\roman*)},nosep,leftmargin=2.2em]
\item the restriction $H^{2}(\cO_{X})\to H^{2}(\cO_{S})$ is injective;
\item no nonzero translation-invariant holomorphic $2$-form on $X$ restricts
  to zero on $S$;
\item $S$ is not Lagrangian for any translation-invariant holomorphic
  symplectic form on $X$;
\item the Gauss image of $S$ in $\operatorname{Gr}(2,V)\subset\PP(\bigwedge^{2}V)$
  spans $\PP(\bigwedge^{2}V)$.
\end{enumerate}
If $e(S)\ne[S]^{2}$, they hold. In particular they hold for every smooth
support of \Cref{tab:smoothsupport}, where $e(S)=12N(36-3N)$ and
$[S]^{2}=24N^{2}$ agree only for $5N=36$.
\end{lemma}

\begin{proof}
The restriction $\iota^{*}\colon H^{2}(X,\CC)\to H^{2}(S,\CC)$ is a morphism
of Hodge structures defined over $\QQ$, so it commutes with complex
conjugation, which interchanges the $(0,2)$ and $(2,0)$ parts; hence
$\iota^{*}$ is injective on $H^{0,2}(X)=H^{2}(\cO_{X})$ if and only if it is
injective on $H^{2,0}(X)=\bigwedge^{2}V^{\vee}$, the translation-invariant
$2$-forms. A holomorphic $2$-form on a compact surface is zero as a class in
$H^{2,0}(S)=H^{0}(K_{S})$ if and only if it is zero as a form; this is
(i)$\Leftrightarrow$(ii). The restriction of $p\wedge q\in\bigwedge^{2}V^{\vee}$
to $S$ is the pull-back of the Pl\"ucker coordinate $p\wedge q$ along the
Gauss map $\gamma\colon S\to\operatorname{Gr}(2,V)$, $s\mapsto T_{s}S$, as a
section of $K_{S}=\gamma^{*}\cO(1)$, so a $2$-form vanishes on $S$ if and
only if the corresponding hyperplane contains $\gamma(S)$; this is
(ii)$\Leftrightarrow$(iv).

A nonzero $2$-form on $V\cong\CC^{4}$ has rank two or four. If $\omega$ has
rank four, $\omega|_{S}=0$ says that every tangent plane of $S$ is
$\omega$-isotropic, that is, $S$ is Lagrangian for $\omega$. If
$\omega=\alpha\wedge\beta$ has rank two, with $\alpha,\beta\in V^{\vee}$
independent, then $\alpha|_{S}$ and $\beta|_{S}$ are independent
holomorphic $1$-forms on $S$, because $S$ generates the simple $X$ and a nonzero
invariant $1$-form restricts nontrivially to a generating subvariety, and
$\alpha|_{S}\wedge\beta|_{S}=0$. By the theorem of Castelnuovo and de
Franchis \cite{BHPV04} there is a fibration
$f\colon S\to C$ onto a curve of genus at least two with $\alpha|_{S}$ and
$\beta|_{S}$ pulled back from $C$. A fibre $F$ of $f$ is a curve in $X$ on
which $\alpha$ and $\beta$ vanish, so all its tangent lines lie in the fixed
plane $B=\ker\alpha\cap\ker\beta\subset V$. The abelian subvariety generated
by $F-F$ is the image of a sum map from a power of $F$, whose differential at
a general point is spanned by tangent vectors of $F$; so its tangent space
lies in $B$, its dimension is at most two, and it is nonzero because $F$ is
not a point. This contradicts the simplicity of $X$. Hence only the
rank-four case occurs, which is (ii)$\Leftrightarrow$(iii).

Finally, if $S$ is Lagrangian for a symplectic $\omega$, then $\omega$
identifies $T_{S}$ with the annihilator of $T_{S}$ in $\Omega^{1}_{X}|_{S}$,
which is the conormal bundle, so $N_{S/X}\cong\Omega^{1}_{S}$ and
$c_{2}(N_{S/X})=e(S)$; since $T_{X}$ is trivial, $c_{2}(N_{S/X})=[S]^{2}$ by
the self-intersection formula. So $e(S)\ne[S]^{2}$ excludes (iii). For the
table, \Cref{thm:smoothinvariants} gives $[S]^{2}=24N^{2}$ and
$e(S)=12N(36-3N)$ at $b=3$, and $24N^{2}=12N(36-3N)$ reads $5N=36$. For
general $b$ the two agree only when $N=\tfrac45b^{2}$, that is
$d=\tfrac35b^{2}$, whose first integral instance is $b=5$, $d=15$.
\end{proof}

\begin{remark}[The class $K_{S}-4\lambda$]\label{rem:gysinkernel}
\Cref{thm:smoothinvariants} records $K_{S}\cdot\lambda=96N$ and
$\lambda^{2}=24N$ at $b=3$, so $K_{S}\cdot\lambda=4\lambda^{2}$, and that
$K_{S}$ is never numerically proportional to $\lambda$. The sharper
statement is that
\[
  \delta=K_{S}-4\lambda\in\ker\bigl(\iota_{*}\colon H^{2}(S,\QQ)\to H^{6}(X,\QQ)\bigr),
  \qquad \delta^{2}=-12N(N-4)<0,
\]
because $\iota_{*}K_{S}=4N\Theta^{3}$ and
$\iota_{*}\lambda=\Theta\cup[S]=N\Theta^{3}$. So $\delta$ is a nonzero
algebraic class of type $(1,1)$, orthogonal by the projection formula to the
whole image of $H^{2}(X,\QQ)$ in $H^{2}(S,\QQ)$, and of negative square: the
Picard number of $S$ is at least two, and the kernel of the Gysin map on
$H^{2}(S,\QQ)$, a sub-Hodge structure, contains an algebraic class. For
general $b$ the class is $K_{S}-\tfrac{4b}{3}\lambda$, with square
$\tfrac{4N}{3}(4b^{2}-9N)$, and its vanishing is the equality case $9N=4b^{2}$
of the Hodge index bound in \Cref{thm:smoothinvariants}. This excludes none
of the four cases; it describes them.
\end{remark}

The third condition is the one that \Cref{thm:factor}(iii)(a) imposes on
every candidate and that the construction of \Cref{setup:markman} meets by
its uniformity over the moduli of principally polarised fourfolds
(\Cref{cor:markmancondition}): the object must extend to first order along
every polarised deformation of $(X,\Theta)$. For $\cF=\cI_{S}(3\Theta)$ with
$S$ smooth this is a condition on $S$. If $S$ extends along a deformation
$X_{v}$ of $X$ that keeps $\Theta$ of type $(1,1)$, then so does $\cF$.
Conversely, the obstruction $v\lrcorner\,\mathrm{At}(\cF)\in\Ext^{2}(\cF,\cF)$
has, under the edge map to $H^{1}(S,\cE xt^{1}(\cF,\cF))=H^{1}(S,N_{S/X})$ of
\Cref{rem:obstructionmoves}, the image of $v$ under
$H^{1}(T_{X})\to H^{1}(\iota^{*}T_{X})\to H^{1}(N_{S/X})$, by
\Cref{lem:edge} applied on the overlaps of a cover; and that image is the
obstruction $\operatorname{ob}_{S}(v)$ to extending the subscheme $S$ along
$X_{v}$. So condition (a) for $\cF$ is the statement that $S$ extends to
first order along every polarised deformation of $X$, and Bloch's
semiregularity map \cite{Blo72} puts it in a form that names the object it
concerns. For $S\subset X$ smooth of codimension two with normal bundle
$N=N_{S/X}$, let
\[
  \pi\colon H^{1}(S,N)\xrightarrow{\ \cong\ }H^{1}(S,N^{\vee}\otimes K_{S})
  \xrightarrow{\ j\ }H^{1}(S,\iota^{*}\Omega^{1}_{X}\otimes K_{S})
  \xrightarrow{\ \iota_{*}\ }H^{3}(X,\Omega^{1}_{X})
\]
be the composite of $N\cong N^{\vee}\otimes\bigwedge^{2}N$, of the conormal
inclusion $N^{\vee}\to\iota^{*}\Omega^{1}_{X}$ tensored with $K_{S}$, and of
the Gysin map, with $K_{X}\cong\cO_{X}$. This is the semiregularity map of
Bloch for codimension two, and it satisfies
$\pi(\operatorname{ob}_{S}(v))=v\lrcorner[S]$ for every $v\in H^{1}(X,T_{X})$.

\begin{lemma}[Duality]\label{lem:blochduality}
Under Serre duality $H^{1}(X,T_{X})\times H^{3}(X,\Omega^{1}_{X})\to H^{4}(\cO_{X})=\CC$,
\[
  \{v\in H^{1}(X,T_{X}):S\ \text{extends to first order along }v\}
  =\pi\bigl(H^{1}(S,N)\bigr)^{\perp}.
\]
\end{lemma}

\begin{proof}
$S$ extends along $v$ if and only if $\operatorname{ob}_{S}(v)=0$, and
$\operatorname{ob}_{S}$ is the composite of the restriction
$r\colon H^{1}(X,T_{X})\to H^{1}(S,\iota^{*}T_{X})$ with the map induced by
$\iota^{*}T_{X}\to N$. The Serre dual of $r$ is the Gysin map
$\iota_{*}\colon H^{1}(S,\iota^{*}\Omega^{1}_{X}\otimes K_{S})\to H^{3}(X,\Omega^{1}_{X}\otimes K_{X})$,
and the Serre dual of $H^{1}(\iota^{*}T_{X})\to H^{1}(N)$ is the map
$H^{1}(N^{\vee}\otimes K_{S})\to H^{1}(\iota^{*}\Omega^{1}_{X}\otimes K_{S})$
induced by the conormal inclusion. So the dual of $\operatorname{ob}_{S}$ is
$\iota_{*}\circ j$, which is $\pi$ after the identification
$N\cong N^{\vee}\otimes K_{S}$, and the kernel of a linear map is the
annihilator of the image of its dual.
\end{proof}

\begin{proposition}[The first-order condition is a rank condition]
\label{prop:smoothrank}
Let $S\subset X$ be a smooth surface with $[S]=N\Theta^{2}$, $N\ne0$. Then
$\pi(H^{1}(S,N))\supseteq\Theta\cup H^{0,2}(X)$, a subspace of dimension six,
and the following are equivalent:
\begin{enumerate}[label=\textup{(\roman*)},nosep,leftmargin=2.2em]
\item $S$ extends to first order along every polarised deformation of
  $(X,\Theta)$, that is, $\cI_{S}(3\Theta)$ satisfies condition (a) of
  \Cref{thm:factor}(iii);
\item $\pi(H^{1}(S,N))=\Theta\cup H^{0,2}(X)$;
\item $\operatorname{rank}\pi=6$, the least value possible.
\end{enumerate}
In particular a smooth secant support is as far from semiregular in Bloch's
sense as a surface with this class can be: $\dim\ker\pi=h^{1}(N)-6$, with
$h^{1}(N)=2h^{0}(N)-2\chi(\cO_{S})+24N^{2}\ge8-2\chi(\cO_{S})+24N^{2}$, which
is at least $88$, $296$, $568$ and $904$ for $N=5,6,7,8$.
\end{proposition}

\begin{proof}
The polarised deformations form $P=\{v:v\lrcorner\Theta=0\}$, of dimension
ten in $H^{1}(T_{X})=V\otimes H^{0,1}(X)$. For $z\in H^{0,2}(X)$ the pairing
is $\langle v,\Theta\cup z\rangle=(v\lrcorner\Theta)\cup z$, so
$\Theta\cup H^{0,2}(X)\subseteq P^{\perp}$; the map $z\mapsto\Theta\cup z$ is
injective by hard Lefschetz, so this subspace has dimension six, and
$\dim P^{\perp}=16-10=6$, so $P^{\perp}=\Theta\cup H^{0,2}(X)$. By
\Cref{lem:blochduality}, (i) says $P\subseteq\pi(H^{1}(N))^{\perp}$, that is
$\pi(H^{1}(N))\subseteq P^{\perp}$. On the other hand
$\pi(\operatorname{ob}_{S}(v))=v\lrcorner[S]=2N\,\Theta\cup(v\lrcorner\Theta)$,
and $v\mapsto v\lrcorner\Theta$ is surjective onto $H^{0,2}(X)$, because
$\Theta$ identifies $V$ with $\overline{V}^{\vee}$ and the contraction becomes
the exterior product $V\otimes\overline{V}^{\vee}\to\bigwedge^{2}\overline{V}^{\vee}$.
So $\pi(H^{1}(N))\supseteq\Theta\cup H^{0,2}(X)$ always, and (i), (ii), (iii)
are equivalent. For $h^{1}(N)$, Riemann--Roch with $c_{1}(N)=K_{S}$ and
$c_{2}(N)=[S]^{2}=24N^{2}$ gives $\chi(N)=2\chi(\cO_{S})-24N^{2}$, Serre
duality with $N^{\vee}\otimes K_{S}\cong N$ gives $h^{2}(N)=h^{0}(N)$, and
$h^{0}(N)\ge4$ because the four translation vector fields map injectively to
$H^{0}(N)$: a translation vector field tangent to $S$ everywhere would restrict
to a nowhere vanishing vector field on $S$, so that $e(S)=c_{2}(T_{S})=0$,
while a secant support has $e(S)=12N(36-3N)\ne0$ for $N=5,6,7,8$ by
\Cref{thm:smoothinvariants}.
\end{proof}

\begin{remark}[What the rank condition asks]\label{rem:smoothrank}
\Cref{prop:smoothrank} does not decide the condition; it converts it into the
statement that the kernel of Bloch's map has codimension exactly six, as it
has for the complete intersections $\Theta\cap\Theta_{x}$, which do move with
$X$. The kernel of $\pi$ is
$j^{-1}\bigl(V^{\vee}\otimes\ker(H^{1}(K_{S})\to H^{3}(\cO_{X}))\bigr)$, and
for $q(S)=4$ it is the image in $H^{1}(N^{\vee}\otimes K_{S})$ of
$H^{0}(\Omega^{1}_{S}\otimes K_{S})$ modulo $V^{\vee}\otimes H^{0}(K_{S})$. A
count with $\chi(\Omega^{1}_{S}\otimes K_{S})=2\chi(\cO_{S})+24N^{2}$ and
$h^{1}(T_{S})\le h^{0}(N)+12$ shows the rank condition to be consistent with
every inequality available, so no exclusion follows at this level; what
would decide it is $\ker\pi$ for an actual surface, and no candidate surface
is known.
\end{remark}

\subsection{Depth and the local obstruction}
\label{ssec:depth}

By \Cref{prop:smoothlocal}, at a smooth point of the support there is
nothing local to compare. This is a special case of a much more general
statement, and the general statement shows exactly which supports
\Cref{thm:candidatefails} and \Cref{cor:localobstruction} exclude and which
lie beyond their reach. The dividing line is neither smoothness nor the local
complete intersection property, but depth.

\begin{theorem}[Cohen--Macaulay supports carry no local obstruction]%
\label{thm:cmvanishing}
Let $Y$ be a smooth variety with $\omega_{Y}\cong\cO_{Y}$, let $Z\subset Y$ be
a closed subscheme of codimension two, and let $p\in Z$. If the local ring
$\cO_{Z,p}$ is Cohen--Macaulay, then
\[
  \cE xt^{q}(\cI_{Z},\cI_{Z})_{p}=0\qquad\text{for every }q\ge2 ,
\]
and the same holds for $\cI_{Z}\otimes L$ for every line bundle $L$.
Consequently both sides of \eqref{eq:heisenberg} have vanishing germ at $p$,
and \Cref{lem:edge} decides nothing there.
\end{theorem}

\begin{proof}
Twisting by $L$ changes no $\cE xt$ sheaf, so take $F=\cI_{Z}$. Write
$A=\cO_{Y,p}$, a regular local ring of dimension $m$, and $I=\cI_{Z,p}$. By
hypothesis $A/I$ is Cohen--Macaulay of dimension $m-2$, so
$\operatorname{depth}A/I=m-2$, and the Auslander--Buchsbaum formula gives
\[
  \operatorname{pd}_{A}(A/I)=m-\operatorname{depth}_{A}(A/I)=2 .
\]
From the exact sequence $0\to I\to A\to A/I\to0$ with $A$ free it follows that
$\operatorname{pd}_{A}I=1$. A module of projective dimension one has a free
resolution of length one, so $\Ext^{q}_{A}(I,M)=0$ for every $A$-module $M$
and every $q\ge2$; taking $M=I$ gives the vanishing, since
$\cE xt^{q}(\cI_{Z},\cI_{Z})_{p}=\Ext^{q}_{A}(I,I)$. The last assertion
follows from \Cref{lem:edge}: the right side of \eqref{eq:heisenberg} is
$z\cdot\id_{F}$ with $z\in H^{2}(\cO_{Y})$ and has vanishing germs, and the
left side lies in a sheaf whose $\cE xt^{2}$ factor is zero.
\end{proof}

\begin{remark}[Sharpness of the hypothesis]%
\label{rem:cmsharp}
Two observations show that \Cref{thm:cmvanishing} is the right statement.

First, Cohen--Macaulayness is strictly weaker than the more familiar
conditions that could be imposed instead. Among codimension two germs in
$\CC^{4}$, the smooth germ $(x_{1},x_{2})$, two planes meeting along a line,
$(x_{1},x_{2}x_{3})$, three planes meeting along a line,
$(x_{1},x_{2}x_{3}x_{4})$, the double plane $(x_{1}^{2},x_{2})$ and the
quadric cone $(x_{1},x_{2}^{2}-x_{3}x_{4})$ are complete intersections. The
fat plane $(x_{1},x_{2})^{2}$ and the cone over the twisted cubic are the
ideals of $2\times2$ minors of
$\left(\begin{smallmatrix}x_{1}&x_{2}&0\\
0&x_{1}&x_{2}\end{smallmatrix}\right)$ and of
$\left(\begin{smallmatrix}x_{1}&x_{2}&x_{3}\\
x_{2}&x_{3}&x_{4}\end{smallmatrix}\right)$, of the expected codimension two,
so they are Cohen--Macaulay by the Hilbert--Burch theorem, although neither
is a local complete intersection, both needing three generators. By
\Cref{thm:cmvanishing}, $\cE xt^{2}(\cI_{Z},\cI_{Z})$ vanishes at all seven.
Reducedness, irreducibility and the complete intersection property are
therefore all irrelevant; what matters is the projective dimension, which by
Auslander--Buchsbaum is a matter of depth.

Second, the vanishing fails at the germs that the construction of
\Cref{setup:markman} produces. Two planes meeting at a point and three planes
meeting pairwise at a point are reduced surface germs that become
disconnected when the point is removed, so they have depth $1$: at least $1$
because they are reduced, and at most $1$ by the connectedness theorem
recalled below. A plane with an embedded point has depth $0$. By Auslander--Buchsbaum the projective
dimensions of $\cO_{Z,p}$ are $3$, $3$ and $4$, and the proof of
\Cref{cor:cmdichotomy}(i) below gives $\cE xt^{q}(\cI_{Z},\cI_{Z})_{p}\ne0$
for $q=2$, $2$ and $3$. By Hartshorne's connectedness theorem, a scheme whose local ring at $p$
is Cohen--Macaulay is connected in codimension one at $p$, so a support with
two components through $p$ meeting only at $p$ is never Cohen--Macaulay
there. This is precisely the configuration analysed in
\Cref{lem:localproducts}.
\end{remark}

\begin{corollary}[The local dichotomy]%
\label{cor:cmdichotomy}
Let $(X,t)$ be a polarised abelian variety of dimension $n\ge4$ and $F$ a
perfect complex with $\ch(F)=au_{t}+bv_{t}$, $(a,b)\ne(0,0)$, whose germ at
each point is the germ of $\cI_{Z}\otimes L$ for a subscheme $Z$ of
codimension two, and let $p\in Z$.
\begin{enumerate}[label=\textup{(\roman*)},leftmargin=2.6em]
\item If $Z$ is not Cohen--Macaulay at $p$, then
  $\cE xt^{q}(\cI_{Z},\cI_{Z})_{p}\ne0$ for $q=\operatorname{pd}\cI_{Z,p}\ge2$,
  so the comparison of germs of \Cref{lem:edge} has room at $p$. If moreover
  $Z$ has at $p$ two smooth branches of codimension two meeting
  transversally, as in \Cref{lem:localproducts} and
  \Cref{cor:localobstruction}, then $F$ does not satisfy the weakened
  criterion, nor does $\Phi(F\boxtimes F_{2})$ for any secant $F_{2}$.
\item If $Z$ is Cohen--Macaulay, then there is no local obstruction at any
  point. Condition (b) of \Cref{thm:factor}(iii) is then a single global
  condition: it holds if and only if the image of $A_{\theta}A_{\theta'}$ in
  $H^{1}\bigl(Z,\cE xt^{1}(\cI_{Z},\cI_{Z})\bigr)$ vanishes for all
  $\theta,\theta'\in H^{0}(T_{X})$, the scalar on the right of
  \eqref{eq:heisenberg} being forced by the trace. Condition (a), the
  extension of $F$ to first order along the polarised deformations, is
  not affected by the depth and remains to be checked.
\end{enumerate}
In case (ii), $\cI_{Z}$ has projective dimension one at every point, so it is
locally the ideal of maximal minors of a Hilbert--Burch matrix; if $Z$ is
moreover smooth, then $\cE xt^{1}(\cI_{Z},\cI_{Z})=N_{Z/X}$ and the invariants
of $Z$ are those of \Cref{thm:smoothinvariants}.
\end{corollary}

\begin{proof}
(i) Put $A=\cO_{X,p}$, of dimension $n$, and $I=\cI_{Z,p}$. Since $A/I$ is not
Cohen--Macaulay, $\operatorname{depth}A/I<\dim A/I\le n-2$, so
$\operatorname{pd}_{A}(A/I)\ge3$ by the Auslander--Buchsbaum formula and
$q=\operatorname{pd}_{A}I\ge2$. In a minimal free resolution
$F_{\bullet}\to I$ the entries of the differentials lie in the maximal ideal
$\mathfrak m$, so $\Ext^{q}_{A}(I,I)$, the cokernel of
$\Hom(F_{q-1},I)\to\Hom(F_{q},I)$, surjects onto
$\Hom(F_{q},I)\otimes A/\mathfrak m\ne0$ by Nakayama's lemma. At two
transversal branches the conclusion is \Cref{thm:candidatefails} and
\Cref{cor:localobstruction}. (ii) \Cref{thm:cmvanishing} gives
$\cE xt^{q}=0$ for $q\ge2$ at every point, so the local-to-global spectral
sequence for $\Ext^{2}(\cI_{Z},\cI_{Z})$ has only the two rows
$H^{1}(\cE xt^{1})$ and $H^{2}(\cO_{X})$, and these lie in different steps of
the filtration, the second being the subgroup of scalars. The right side of
\eqref{eq:heisenberg} lies in that subgroup, so its image in the quotient
$H^{1}(\cE xt^{1})$ is zero, and (b) forces the image of the left side there
to vanish. Conversely, if that image vanishes, then
$\mathrm{At}(F)^{2}=z\cdot\id_{F}$ for some
$z\in H^{2}(\cO_{X})\otimes\bigwedge^{2}H^{1,0}$, the map
$z\mapsto z\cdot\id_{F}$ being injective because its composite with the
trace is multiplication by $\rk F=1$; taking traces,
$z=\operatorname{tr}\mathrm{At}(F)^{2}=-d\,t^{2}$ by the last assertion of
\Cref{thm:factor}(iii), which is (b). The statement on projective dimension
was established in the proof of \Cref{thm:cmvanishing}, and the Hilbert--Burch theorem is the
structure theorem for ideals of projective dimension one over a regular local
ring; the last clause is \Cref{prop:smoothlocal} together with
\Cref{thm:smoothinvariants}.
\end{proof}

\begin{remark}[Normalisation and depth]%
\label{rem:normalisationcm}
\Cref{cor:cmdichotomy} explains the tension recorded in
\Cref{rem:normalisedpoints}. The partial normalisation performed in
\Cref{setup:markman} separates the two branches at an isolated double point,
and separating them is exactly what would make the local ring
Cohen--Macaulay there; the normalisation is thus an attempt to reach case
(ii). \Cref{rem:normalisedpoints} shows that the attempt cannot succeed:
normalising every isolated point moves $\chi(\cO_{\widetilde Z})$ to $50N$,
while the secant plane demands $(d+9)(27-d)=72N-4N^{2}$, and the two agree
only at $N=11/2$. Repairing the depth destroys the Chern character, and
keeping the Chern character keeps a point where the depth fails. This is the
precise form of the obstruction, and it explains why a construction
assembled from intersections of translates of a polarisation cannot be made
to work: by \Cref{cor:localobstruction} such a chain always has a point of
the bad type, and by \Cref{rem:normalisedpoints} it cannot be separated
without leaving the plane.
\end{remark}

\subsection{Hilbert--Burch resolutions}
\label{ssec:hilbertburch}

\Cref{cor:cmdichotomy} leaves one class of supports and, beside the
extension along the polarised deformations, one condition on them.
The class consists of the Cohen--Macaulay subschemes of codimension two, and
for these the proof of \Cref{thm:cmvanishing} gives more than the vanishing
it was used for: the ideal sheaf has projective dimension one at every point.
This is the Hilbert--Burch shape, whose invariants are polynomials in the twists.

\begin{setupenv}[A resolution of the support]\label{setup:burch}
Let $(X,\Theta)$ be a principally polarised abelian fourfold with
$\NS(X)_{\QQ}=\QQ\Theta$, let $Z\subset X$ be Cohen--Macaulay of codimension
two, and suppose $\cI_{Z}$ admits a global resolution
\begin{equation}\label{eq:burchres}
  0\longrightarrow E_{1}\stackrel{\varphi}{\longrightarrow}E_{0}
  \longrightarrow\cI_{Z}\longrightarrow0,
  \qquad \rk E_{0}=r+1,\quad\rk E_{1}=r,
\end{equation}
with $E_{0}=\bigoplus_{i=1}^{r+1}\cO(e_{i}\Theta)$ and
$E_{1}=\bigoplus_{j=1}^{r}\cO(f_{j}\Theta)$ sums of line bundles. Put
$p_{i}=e_{i}+b$ and $q_{j}=f_{j}+b$.
\end{setupenv}

\begin{lemma}[The resolution in moments]\label{lem:burchmoments}
In \Cref{setup:burch}, $\ch(\cI_{Z}(b\Theta))=u_{\Theta}+b\,v_{\Theta}$ holds
if and only if
\begin{equation}\label{eq:burchsystem}
  \sum_{i=1}^{r+1}p_{i}^{\,k}-\sum_{j=1}^{r}q_{j}^{\,k}=c_{k},
  \qquad k=0,1,2,3,4,
  \qquad (c_{k})=(1,\;b,\;-d,\;-db,\;d^{2}).
\end{equation}
The equation at $k=0$ is the rank count and is automatic. In addition,
$\varphi$ can be injective only if $\max_{i}e_{i}\ge\max_{j}f_{j}$, and
$\cO_{X}$ embeds in $E_{0}^{\vee}$ only if $\min_{i}e_{i}\le0$.
\end{lemma}

\begin{proof}
From \eqref{eq:burchres}, $\ch(\cI_{Z})=\ch(E_{0})-\ch(E_{1})$, so
$\ch(\cI_{Z}(b\Theta))=\sum_{i}e^{p_{i}\Theta}-\sum_{j}e^{q_{j}\Theta}$.
Comparing coefficients of $\Theta^{k}/k!$ with those of
$u_{\Theta}+bv_{\Theta}$, which are the $c_{k}$ displayed, gives
\eqref{eq:burchsystem}. For the last two statements, recall that on an
abelian variety $\Hom(\cO(f\Theta),\cO(e\Theta))=H^{0}(\cO((e-f)\Theta))$
vanishes unless $e\ge f$. If $\varphi$ is injective, each summand of $E_{1}$
must map nonzero to $E_{0}$, which requires $\max_{i}e_{i}\ge f_{j}$ for
every $j$. Dualising \eqref{eq:burchres} and using $\cI_{Z}^{\vee}=\cO_{X}$
gives an injection $\cO_{X}\hookrightarrow E_{0}^{\vee}$, which requires
$-e_{i}\ge0$ for some $i$.
\end{proof}

The case $r=1$ consists of the complete intersections, and it is settled for
every $b$ and every $d$ by an argument that needs no search.

\begin{theorem}[No complete intersection supports a secant object]%
\label{thm:noci}
Let $(X,\Theta)$ be a principally polarised abelian fourfold with
$\NS(X)_{\QQ}=\QQ\Theta$, let $d>0$ and $b\ge1$. No complete intersection
$Z=D_{1}\cap D_{2}$ of two divisors satisfies
$\ch(\cI_{Z}(b\Theta))=u_{\Theta}+b\,v_{\Theta}$.
\end{theorem}

\begin{proof}
Such a $Z$ is \Cref{setup:burch} with $r=1$, the resolution being the Koszul
complex of the two equations. Put $w(x)=x^{2}+d$, which is strictly positive
on $\RR$ because $d>0$, and read \eqref{eq:burchsystem} in the combinations
$c_{k+2}+d\,c_{k}$ for $k=0,1,2$. Those three numbers are
\[
  c_{2}+dc_{0}=-d+d=0,\qquad
  c_{3}+dc_{1}=-db+db=0,\qquad
  c_{4}+dc_{2}=d^{2}-d^{2}=0,
\]
so that
\[
  \sum_{i}w(p_{i})\,p_{i}^{\,k}=\sum_{j}w(q_{j})\,q_{j}^{\,k},
  \qquad k=0,1,2 .
\]
The two sides are the moments of order $0$, $1$ and $2$ of the positive
measures $\nu_{P}=\sum_{i}w(p_{i})\delta_{p_{i}}$ and
$\nu_{Q}=\sum_{j}w(q_{j})\delta_{q_{j}}$, so those measures have the same
mass, the same mean and the same variance. At $r=1$ the measure $\nu_{Q}$ has
one atom, hence variance zero, so $\nu_{P}$ has variance zero as well, and a
positive measure with two atoms has variance zero only if the atoms coincide.
Hence $p_{1}=p_{2}$, and equality of the means forces $q_{1}=p_{1}$. The
resolution \eqref{eq:burchres} then cancels, and $\cI_{Z}\cong\cO(e_{1}\Theta)$
would be a line bundle, which the ideal sheaf of a subscheme of codimension
two is not.

In terms of the divisors, the same computation shows that their degrees
$a_{1}$ and $a_{2}$ would have to satisfy $a_{1}+a_{2}=\tfrac43b$ and
$(b-a_{1})(b-a_{2})=\tfrac12\bigl(\tfrac13b^{2}+d\bigr)$; the quadratic with
roots $b-a_{1}$ and $b-a_{2}$ then has discriminant $-\tfrac29b^{2}-2d<0$, so
the degrees are not even real.
\end{proof}

For $r\ge2$ the moment conditions alone no longer decide the question. The
Chern characters of the bundles of \Cref{setup:burch}, and of every bundle
built from line bundles, lie in one lattice, and the secant plane meets that
lattice only along a congruence. Write a class $\xi\in\QQ[\Theta]/(\Theta^{5})$
as $\xi=\sum_{k=0}^{4}c_{k}(\xi)\,\Theta^{k}/k!$, as in
\Cref{lem:burchmoments}, so that $c_{k}(e^{j\Theta})=j^{k}$, and let
$\cL_{\Theta}$ be the subgroup generated by the classes
$e^{j\Theta}=\ch\cO_{X}(j\Theta)$, $j\in\ZZ$. It is a subring, since
$e^{a\Theta}e^{b\Theta}=e^{(a+b)\Theta}$, and it contains $\ch$ of every
bundle obtained from line bundles algebraically equivalent to multiples of
$\Theta$ by direct sums, tensor products, duals, exterior and symmetric
powers, kernels of surjections and cokernels of injections between such
bundles, and extensions;
in particular the kernel bundles
$M_{m\Theta}=\ker\bigl(H^{0}(\cO(m\Theta))\otimes\cO_{X}\to\cO_{X}(m\Theta)\bigr)$,
of class $m^{4}-e^{m\Theta}$, and their duals and twists.

\begin{lemma}[Binomial moments]\label{lem:binomialmoments}
Define $m_{k}(\xi)\in\QQ$, $0\le k\le4$, by
\begin{equation}\label{eq:stirling}
  c_{i}(\xi)=\sum_{k=0}^{i}S(i,k)\,k!\,m_{k}(\xi),
\end{equation}
where $S(i,k)$ are the Stirling numbers of the second kind; explicitly
\begin{gather*}
  m_{0}=c_{0},\qquad m_{1}=c_{1},\qquad m_{2}=\tfrac12(c_{2}-c_{1}),\\
  m_{3}=\tfrac16(c_{3}-3c_{2}+2c_{1}),\qquad
  m_{4}=\tfrac1{24}(c_{4}-6c_{3}+11c_{2}-6c_{1}).
\end{gather*}
Then $\xi\in\cL_{\Theta}$ if and only if $m_{k}(\xi)\in\ZZ$ for $k=0,\dots,4$.
\end{lemma}

\begin{proof}
For $\xi=\sum_{j}a_{j}e^{j\Theta}$ one has $c_{i}(\xi)=\sum_{j}a_{j}j^{i}$,
and $j^{i}=\sum_{k}S(i,k)\,k!\binom{j}{k}$, so \eqref{eq:stirling} holds
with $m_{k}(\xi)=\sum_{j}a_{j}\binom{j}{k}\in\ZZ$. Conversely, the five
vectors $\bigl(\binom{j}{0},\dots,\binom{j}{4}\bigr)$ for $j=0,\dots,4$
form a unitriangular matrix, so every integer vector $(m_{0},\dots,m_{4})$
is an integer combination of them, and the corresponding
$\sum_{j=0}^{4}a_{j}e^{j\Theta}$ lies in $\cL_{\Theta}$ with the prescribed
moments. Since \eqref{eq:stirling} is invertible over $\QQ$, the moments
determine $\xi$.
\end{proof}

Equivalently, sending $e^{j\Theta}$ to $x^{j}$ identifies
$\QQ[\Theta]/(\Theta^{5})$ with $\QQ[y]/(y^{5})$, $y=x-1$, the class
$e^{j\Theta}$ going to $(1+y)^{j}$; then $m_{k}(\xi)$ is the coefficient of
$y^{k}$ and $\cL_{\Theta}$ is $\ZZ[y]/(y^{5})$.

\begin{theorem}[The secant plane and the lattice of line bundles]
\label{thm:lattice}
For $(a,b)\in\ZZ^{2}$ and $d\ge1$,
\[
  m\bigl(a\,u_{\Theta}+b\,v_{\Theta}\bigr)
  =\Bigl(a,\ b,\ -\frac{ad+b}{2},\ \frac{3ad+2b-bd}{6},\
  \frac{ad^{2}-11ad+6bd-6b}{24}\Bigr).
\]
Consequently $a\,u_{\Theta}+b\,v_{\Theta}\in\cL_{\Theta}$ if and only if
\begin{gather*}
  ad\equiv b\pmod 2,\qquad 3ad+2b-bd\equiv0\pmod 6,\\
  ad^{2}-11ad+6bd-6b\equiv0\pmod{24}.
\end{gather*}
For $a=1$ and $b=3$ the moments are
$\bigl(1,3,-\tfrac{d+3}{2},1,\tfrac{(d+9)(d-2)}{24}\bigr)$, and
$u_{\Theta}+3v_{\Theta}\in\cL_{\Theta}$ if and only if $d$ is odd and
$24\mid(d+9)(d-2)$, that is, if and only if
\[
  d\equiv15\ \text{or}\ 23\pmod{24}.
\]
\end{theorem}

\begin{proof}
In the variable $y=x-1$ put
$L=\log(1+y)=y-\tfrac12y^{2}+\tfrac13y^{3}-\tfrac14y^{4}$ modulo $y^{5}$.
Under $e^{j\Theta}\mapsto x^{j}$ the class $\Theta$ goes to $L$, hence
$u_{\Theta}\mapsto1-\tfrac d2L^{2}+\tfrac{d^{2}}{24}L^{4}$ and
$v_{\Theta}\mapsto L-\tfrac d6L^{3}$. With
$L^{2}=y^{2}-y^{3}+\tfrac{11}{12}y^{4}$, $L^{3}=y^{3}-\tfrac32y^{4}$ and
$L^{4}=y^{4}$ modulo $y^{5}$, the coefficients of $1,y,\dots,y^{4}$ in
$a\,u_{\Theta}+b\,v_{\Theta}$ are the five displayed expressions, and the
congruences are the integrality of the last three. For $b=3$ the second
reads $3d+6-3d=6$, so $m_{3}=1$, and the third reads
$d^{2}+7d-18=(d+9)(d-2)$. Modulo $8$ the condition is $8\mid(d+1)(d-2)$
with $d$ odd, hence $d\equiv7\pmod8$; modulo $3$ it is $3\mid d(d-2)$, hence
$d\equiv0,2\pmod3$; together $d\equiv15,23\pmod{24}$.
\end{proof}

\begin{corollary}[Split resolutions, for every rank and every twist]
\label{cor:latticeburch}
Let $Z\subset X$ be Cohen--Macaulay of codimension two with a resolution
$0\to E_{1}\to E_{0}\to\cI_{Z}\to0$ in which $\ch E_{0}$ and $\ch E_{1}$ both
lie in $\cL_{\Theta}$: for instance \Cref{setup:burch} with any rank and any
twists, or $E_{0}$ and $E_{1}$ built from such line bundles and from the
bundles $M_{m\Theta}$ by the operations listed above. If
$\ch(\cI_{Z}(3\Theta))=u_{\Theta}+3v_{\Theta}$, then $d\equiv15$ or
$23\pmod{24}$. In particular this never happens for the four discriminants
of \Cref{cor:fourdiscriminants}.
\end{corollary}

\begin{proof}
$\ch(\cI_{Z}(3\Theta))=e^{3\Theta}(\ch E_{0}-\ch E_{1})\in\cL_{\Theta}$, and
\Cref{thm:lattice} applies.
\end{proof}

\begin{remark}[What the congruence measures]\label{rem:latticedefect}
The fractional part of $m_{4}$ is the invariant that separates the secant
plane from the line bundles. For $d\in\{1,3,5,7\}$ it is
$-\tfrac{5}{12},\tfrac12,\tfrac34,\tfrac13$ modulo $\ZZ$, and any object with
Chern character $u_{\Theta}+3v_{\Theta}$ must carry a class outside
$\cL_{\Theta}$ with exactly that defect. Two checks. A point has
$m=(0,0,0,0,1/24)$, and the partial normalisation of \Cref{setup:markman} at
$(d+9)(2d-1)$ isolated points shifts $m_{4}$ by $(d+9)(2d-1)/24$; since
$(d+9)(d-2)+(d+9)(2d-1)=3(d+9)(d-1)\equiv0\pmod{24}$ for odd $d$, the
unnormalised chain, which lies in $\cL_{\Theta}$, is moved onto the plane by
exactly the right defect, as it must be. For a smooth support the class of
the structure sheaf has $m(\cO_{S})=\bigl(0,0,N,-3N,N(83-2N)/12\bigr)$, and
$N(83-2N)/12$ is not an integer for $N\in\{5,6,7,8\}$; so a smooth secant
support is a surface whose structure sheaf is not in the line-bundle
subring, unlike every complete intersection of translates of $\Theta$ and
every degeneracy locus of a map between bundles in $\cL_{\Theta}$. The
Fourier--Mukai transforms of line bundles and the push-forwards of line
bundles along isogenies do have classes outside $\cL_{\Theta}$, which is why
the semi-homogeneous bundles need the separate argument of
\Cref{thm:noblocks}.
\end{remark}

The lattice constrains the classes of the two terms of a resolution; their
ranks constrain the Chern classes, because a vector bundle of rank $s$ has
$c_{k}=0$ for $k>s$, and on a fourfold this bites for $s\le3$.

\begin{proposition}[Resolutions of rank one and two]\label{prop:lowrankburch}
Let $(X,\Theta)$ be a principally polarised abelian fourfold, let $b\ge1$
and $d>0$, and let $Z\subset X$ be of codimension two with a resolution
\[
  0\longrightarrow E_{1}\longrightarrow E_{0}\longrightarrow\cI_{Z}
  \longrightarrow0
\]
by vector bundles of ranks $r$ and $r+1$, and with
$\ch(\cI_{Z}(b\Theta))=u_{\Theta}+b\,v_{\Theta}$.
\begin{enumerate}[label=\textup{(\roman*)},nosep,leftmargin=2.2em]
\item $r\ge2$, whatever the bundles. Equivalently, $Z$ is not the zero locus
  of a section of a vector bundle of rank two.
\item If $r=2$ and the Chern classes of $E_{1}$ are polynomials in $\Theta$,
  as they are when $\Hdg^{2}(X)=\QQ\Theta$ and $\Hdg^{4}(X)=\QQ\Theta^{2}$,
  then, writing $c_{1}(E_{1}(b\Theta))=a\Theta$ with $a\in\ZZ$,
  \[
    3d-b^{2}-4ab\equiv0\pmod{12}.
  \]
  At $b=3$ this says $d\equiv3\pmod4$.
\end{enumerate}
In particular, at the discriminants $d=1$ and $d=5$ of
\Cref{cor:fourdiscriminants}, a resolution whose Chern classes are
polynomials in $\Theta$ has $r\ge3$.
\end{proposition}

\begin{proof}
Put $E=E_{1}(b\Theta)$, $G=E_{0}(b\Theta)$ and $F=\cI_{Z}(b\Theta)$, so that
$0\to E\to G\to F\to0$ is exact and $c(G)=c(E)\,c(F)$. Newton's identities
applied to $\ch(F)=1+b\Theta-\frac d2\Theta^{2}-\frac{db}{6}\Theta^{3}
+\frac{d^{2}}{24}\Theta^{4}$ give, with $N=\frac12(b^{2}+d)$,
\[
  c(F)=1+b\Theta+N\Theta^{2}+\tfrac{bN}{3}\Theta^{3}
  +\tfrac{N(b^{2}-3d)}{12}\Theta^{4}.
\]
Since $E$ and $G$ are vector bundles, $c_{k}(E)=0$ for $k>r$ and
$c_{k}(G)=0$ for $k>r+1$.

(i) Let $r=1$ and $\ell=c_{1}(E)$, so that $c(G)=(1+\ell)\,c(F)$. The rank
two bundle $G$ has $c_{3}(G)=\frac{bN}{3}\Theta^{3}+N\ell\,\Theta^{2}=0$, and
cup product with $\Theta^{2}$ is an isomorphism $H^{2}(X,\QQ)\to H^{6}(X,\QQ)$
by hard Lefschetz, so $\ell=-\frac b3\Theta$. Then
\[
  c_{4}(G)=c_{4}(F)+\ell\,c_{3}(F)
  =N\Bigl(\frac{b^{2}-3d}{12}-\frac{b^{2}}{9}\Bigr)\Theta^{4}
  =-\frac{N(b^{2}+9d)}{36}\,\Theta^{4},
\]
which is not zero, since $N>0$ and $d>0$. A section of a rank two bundle $V$
vanishing in codimension two exactly on $Z$ gives the Koszul resolution
$0\to\det V^{\vee}\to V^{\vee}\to\cI_{Z}\to0$, which has $r=1$; conversely,
at $r=1$ the map $E_{1}\to E_{0}$ is a section of $E_{0}\otimes E_{1}^{\vee}$
with zero locus $Z$.

(ii) Let $r=2$, and write $c_{1}(E)=a\Theta$ and $c_{2}(E)=\frac m2\Theta^{2}$
with $a,m\in\ZZ$, since $\Theta$ and $\frac12\Theta^{2}$ are primitive
integral classes. Now $c_{3}(E)=c_{4}(E)=0$, and the rank three bundle $G$
has
\[
  c_{4}(G)=c_{4}(F)+c_{1}(E)\,c_{3}(F)+c_{2}(E)\,c_{2}(F)=0 .
\]
Integrating, with $\int_{X}\Theta^{4}=24$, and dividing by $2N$ gives
$6m=3d-b^{2}-4ab$. By Riemann--Roch with trivial Todd class,
$\chi(E)=\int_{X}\ch_{4}(E)=a^{4}-2a^{2}m+\frac12m^{2}$ is an integer, so $m$
is even, which is the congruence. At $b=3$ it reads $3(d-3)\equiv12a\equiv0$
modulo $12$, that is, $d\equiv3\pmod4$.
\end{proof}

\begin{remark}[What the ranks exclude]\label{rem:lowrankburch}
Part (i) extends \Cref{thm:noci} from two divisors to the zero locus of a
section of any rank two bundle, stable or not, so the construction of Serre
produces no secant support; it holds for every $X$, with no condition on
$\NS(X)$. Part (ii) holds for arbitrary bundles whose Chern classes are
polynomials in $\Theta$, not only for sums of line bundles, and at $b=3$ it
excludes $r=2$ whenever $d\not\equiv3\pmod4$. The integrality of
the Chern classes of $E_{0}$ is automatic, since
$c(E_{0}(b\Theta))=c(E_{1}(b\Theta))\,c(F)$ and $c(F)$ is integral, so it is
the vanishing above the rank, and not integrality, that constrains. Neither
condition restricts $r\ge3$: a bundle of rank four has no vanishing Chern
class on a fourfold, and at $r=3$ the only condition, $c_{4}(E)=0$, involves
$E$ alone.
\end{remark}

\Cref{cor:latticeburch} says where a split resolution can have the right
Chern character. A stronger statement holds: no object with a split
resolution satisfies the criterion, at any rank and for any discriminant
(\Cref{thm:nosplit}). It
rests on a necessary condition that we isolate, because it applies to
every object and can be tested on any candidate before any global
computation is attempted.

\begin{lemma}[Morphisms to a line bundle are killed by $H^{2}(\cO)$]
\label{lem:annihilation}
Let $(X,t)$ be a polarised abelian variety of dimension $n\ge3$, let $F$ be
a perfect complex with $\ch(F)=au_{t}+bv_{t}$, $(a,b)\ne(0,0)$, which
satisfies the weakened criterion of \Cref{setup:criterion}, and let $L$ be a
line bundle with $c_{1}(L)\in\QQ\,t$. Then for every $m\in\ZZ$, every
morphism $p\colon F\to L[m]$, every morphism $q\colon L[m]\to F$, and every
$z\in H^{2}(\cO_{X})$,
\[
  z\cdot p=0\ \ \text{in }\Hom(F,L[m+2]),\qquad
  z\cdot q=0\ \ \text{in }\Hom(L[m],F[2]),
\]
where $z\cdot(\,\cdot\,)$ is the $H^{*}(\cO_{X})$-module structure on
morphisms in $D^{b}(X)$. In particular, if $F=\cI_{Z}\otimes L'$ for a closed
subscheme $Z$ and a line bundle $L'$ with $c_{1}(L')\in\QQ t$, then the
restriction $H^{n-2}(\cO_{X})\to H^{n-2}(\cO_{Z})$ is injective.
\end{lemma}

\begin{proof}
For a constant bivector $\pi$ the class
$x_{\pi}=(\tfrac d2\,\pi\lrcorner\,t^{2},0,\pi)$ preserves $\ch(F)$ by
\Cref{thm:factor}(ii), so $\operatorname{ev}_{F}(x_{\pi})=0$ by hypothesis.
The Hochschild action is by natural transformations $\id\to\id[2]$ of
$D^{b}(X)$, so for $p\colon F\to L[m]$ we have
$\operatorname{ev}_{L[m]}(x_{\pi})\circ p=p\circ\operatorname{ev}_{F}(x_{\pi})=0$,
where $\operatorname{ev}_{L[m]}=\operatorname{ev}_{L}[m]$. By
\eqref{eq:evsplit}, $\operatorname{ev}_{L}(x_{\pi})=c_{\pi}\cdot\id_{L}$ with
$c_{\pi}=\tfrac12\,\pi\lrcorner\,(\lambda^{2}+d\,t^{2})\in H^{0,2}(X)$ and
$\lambda=c_{1}(L)=e\,t$, that is,
$c_{\pi}=\tfrac{e^{2}+d}{2}\,\pi\lrcorner\,t^{2}$. Since $e^{2}+d>0$ and
$\pi\mapsto\pi\lrcorner\,t^{2}$ is an isomorphism from
$H^{0}(\bigwedge^{2}T_{X})$ onto $H^{0,2}(X)$, as in the proof of
\Cref{thm:factor}(iii), $c_{\pi}$ runs over all of $H^{2}(\cO_{X})$ as $\pi$
varies; as $(c_{\pi}\cdot\id_{L})\circ p=c_{\pi}\cdot p$, this proves the
assertion for $p$. The assertion for $q$ follows in the same way from the
square
$\operatorname{ev}_{F}(x_{\pi})\circ q=q\circ\operatorname{ev}_{L[m]}(x_{\pi})$.
For the last assertion, apply the first to the inclusion
$p\colon\cI_{Z}\otimes L'\to L'$. After twisting by $L'^{-1}$, the sequence
$0\to\cI_{Z}\to\cO_{X}\to\cO_{Z}\to0$ gives an exact sequence
\[
  \Ext^{2}(\cO_{Z},\cO_{X})\longrightarrow H^{2}(\cO_{X})
  \longrightarrow\Ext^{2}(\cI_{Z},\cO_{X}),
\]
in which the second map is $z\mapsto z\cdot p$. So $z\cdot p=0$ for all $z$
means that the first map is surjective, and by Serre duality, $\omega_{X}$
being trivial, the first map is dual to the restriction
$H^{n-2}(\cO_{X})\to H^{n-2}(\cO_{Z})$.
\end{proof}

\begin{theorem}[No two-term complex of powers of the polarisation satisfies
the criterion]\label{thm:nosplit}
Let $(X,t)$ be a polarised abelian variety of dimension $n\ge3$ and $M$ an
ample line bundle with $c_{1}(M)\in\QQ t$. Let $A$ and $B$ be direct sums of
tensor powers $M^{\otimes e}$, $e\in\ZZ$, let $\varphi\colon A\to B$ be any
morphism, and let $F$ be the cone of $\varphi$, or any shift of it. If $F\ne0$
and $\ch(F)=au_{t}+bv_{t}$ with $(a,b)\ne(0,0)$, then $F$ does not satisfy
the weakened criterion. In particular, no object with a resolution
$0\to E_{1}\to E_{0}\to F\to0$ by sums of powers of $M$ satisfies it,
whatever the rank.
\end{theorem}

\begin{proof}
A pair of summands $M^{\otimes e}$ of $A$ and $M^{\otimes e}$ of $B$ joined
by a nonzero entry of $\varphi$, which is then a nonzero scalar, can be
cancelled without changing the cone up to isomorphism; do this until no such
pair remains. If $A=0$ or $B=0$ after cancelling, $F$ is a shifted sum of
line bundles, $\operatorname{ev}_{F}(x_{\pi})$ is diagonal with the nonzero
entries $c_{\pi}\cdot\id$ of the proof of \Cref{lem:annihilation}, and the
criterion fails as in \Cref{cor:noline}; so suppose both are nonzero and
consider the triangle $A\xrightarrow{\varphi}B\xrightarrow{\;q\;}F\to A[1]$.
Let $M^{\otimes b_{j}}$ be a summand of $B$, $\iota_{j}$ its inclusion into
$B$, and $q_{j}=q\circ\iota_{j}\colon M^{\otimes b_{j}}\to F$. Applying
$\Hom(M^{\otimes b_{j}},-)$ to the triangle gives the exact sequence
\[
  \Hom(M^{\otimes b_{j}},A[2])\xrightarrow{\;\varphi_{*}\;}
  \Hom(M^{\otimes b_{j}},B[2])\xrightarrow{\;q_{*}\;}
  \Hom(M^{\otimes b_{j}},F[2]),
\]
and $z\cdot q_{j}=q_{*}(z\cdot\iota_{j})$ for $z\in H^{2}(\cO_{X})$. Now
$\Hom(M^{\otimes b_{j}},B[2])=\bigoplus_{k}H^{2}(M^{\otimes(b_{k}-b_{j})})$,
and $H^{2}(M^{\otimes m})=0$ for $m\ne0$: for $m>0$ by Kodaira vanishing,
$\omega_{X}$ being trivial, and for $m<0$ because a negative line bundle on
an abelian variety has cohomology only in degree $n\ge3$. So $z\cdot\iota_{j}$
is the element with component $z$ in the summand $k=j$ and zero elsewhere. It
lies in the image of $\varphi_{*}$ only if there is
$\psi\in\Hom(M^{\otimes b_{j}},A[2])=\bigoplus_{i}H^{2}(M^{\otimes(a_{i}-b_{j})})$
with $(\varphi\circ\psi)_{j}=\sum_{i}\varphi_{ji}\psi_{i}=z$; but $\psi_{i}\ne0$
forces $a_{i}=b_{j}$, and for such $i$ the entry $\varphi_{ji}$ is a scalar,
zero after the cancellation. So $z\cdot q_{j}\ne0$ for every $z\ne0$, which
contradicts \Cref{lem:annihilation}. The last assertion is the case of a
resolution, where $F\cong\coker\varphi$ is the cone of $\varphi\colon E_{1}\to E_{0}$.
\end{proof}

The powers of the polarisation are not the only bundles on which the
compensated Poisson classes act by scalars. The same holds for every bundle
whose Atiyah class is a multiple of the identity, and on an abelian variety
these are the semi-homogeneous bundles of Mukai \cite{Muk78}. The argument of
\Cref{thm:nosplit} goes through for them, with the cohomology of the fibre
product of two isogenies in place of the cohomology of a power of $\Theta$.
We state the result in a self-contained form.

\begin{definition}[Isogeny blocks]\label{def:block}
Let $(X,\Theta)$ be a polarised abelian variety. An \emph{isogeny block} is a
vector bundle $E=\pi_{*}L$, where $\pi\colon Y\to X$ is an isogeny and $L$ is
a line bundle on $Y$ with $c_{1}(L)=\nu\,\pi^{*}\Theta$ in $H^{2}(Y,\QQ)$,
$\nu\in\QQ$; its \emph{slope} is $\nu$, so that $c_{1}(E)=\nu\deg(\pi)\,\Theta$.
The block is \emph{simple} if $\tau_{g}^{*}L\not\cong L$ for every nonzero
$g\in\ker\pi$. The powers $\cO(e\Theta)$ are the blocks with $\pi=\id$.
\end{definition}

\begin{lemma}[Three properties of a block]\label{lem:blocks}
Let $E=\pi_{*}L$ be an isogeny block of slope $\nu$ on an abelian $n$-fold.
\begin{enumerate}[label=\textup{(\roman*)},nosep]
\item $\mathrm{At}(E)=\nu\Theta\cdot\id_{E}$; hence for a constant bivector
  $\pi$ the compensated class $x_{\pi}$ of \Cref{thm:factor}(ii) acts by
  $\operatorname{ev}_{E}(x_{\pi})=\tfrac{\nu^{2}+d}{2}\,(\pi\lrcorner\,\Theta^{2})\cdot\id_{E}$,
  and \Cref{lem:annihilation} holds with $E$ in place of the line bundle $L$.
\item For a second block $E'=\pi'_{*}L'$ of slope $\nu'$,
  $\Ext^{q}(E',E)=H^{q}(Z,M)$ with $Z=Y'\times_{X}Y$, a disjoint union of
  abelian varieties isogenous to $X$, and $M$ a line bundle whose class on
  each component is $(\nu-\nu')$ times the pullback of $\Theta$. If
  $\nu\ne\nu'$ and $n\ge3$ then $\Ext^{2}(E',E)=0$. If $E$ is simple then
  $\Ext^{2}(E,E)=H^{2}(\cO_{X})\cdot\id_{E}$.
\item A simple block is $\mu$-stable with respect to $\Theta$; hence a
  nonzero morphism between simple blocks of the same slope is an
  isomorphism.
\end{enumerate}
\end{lemma}

\begin{proof}
(i) Since $\pi$ is \'etale, $\pi^{*}\Omega^{1}_{X}=\Omega^{1}_{Y}$ and $\pi_{*}$
is exact, so the first jet sequence of $\pi_{*}L$ is the direct image of
that of $L$, and $\mathrm{At}(\pi_{*}L)$ is the image of $\mathrm{At}(L)$
under $\pi_{*}\colon\Ext^{1}_{Y}(L,L\otimes\pi^{*}\Omega^{1}_{X})\to
\Ext^{1}_{X}(\pi_{*}L,\pi_{*}L\otimes\Omega^{1}_{X})$. As $\mathrm{At}(L)=c_{1}(L)\cdot\id_{L}
=\nu\,\pi^{*}\Theta\cdot\id_{L}$, the projection formula gives
$\nu\Theta\cdot\id_{\pi_{*}L}$. Then \eqref{eq:evgeneral} with a scalar Atiyah
class gives the value of $\operatorname{ev}_{E}(x_{\pi})$, and the proof of
\Cref{lem:annihilation} used nothing else about $L$.

(ii) Adjunction for the finite \'etale $\pi'$, flat base change along the
cartesian square with vertices $Z$, $Y'$, $Y$, $X$, and the projection formula
give $\Ext^{q}_{X}(\pi'_{*}L',\pi_{*}L)=H^{q}(Y',L'^{-1}\otimes\pi'^{*}\pi_{*}L)
=H^{q}(Z,p'^{*}L'^{-1}\otimes p^{*}L)$, with $p,p'$ the two projections; the
class of $M=p'^{*}L'^{-1}\otimes p^{*}L$ is $(\nu-\nu')$ times the pullback of
$\Theta$ on every component, which is an abelian variety finite \'etale over
$Y'$. A line bundle on an abelian $n$-fold whose class is a positive
multiple of an ample class has cohomology only in degree $0$, and one whose
class is a negative multiple only in degree $n$, by Kodaira vanishing and
Serre duality with trivial canonical bundle; so $H^{2}(Z,M)=0$ when
$\nu\ne\nu'$ and $n\ge3$. When $E'=E$, the components of $Y\times_{X}Y$ are
the graphs of the translations by $g\in\ker\pi$, on which $M$ restricts to
$L^{-1}\otimes\tau_{g}^{*}L\in\Pic^{0}(Y)$; a nontrivial element of $\Pic^{0}$
has no cohomology, so $\Ext^{q}(E,E)=\bigoplus_{g}H^{q}(\cO_{Y})$ over the
$g$ with $\tau_{g}^{*}L\cong L$, which for a simple block is $g=0$ alone, and
$H^{q}(\cO_{Y})=H^{q}(\cO_{X})$ through $\pi$.

(iii) The pullback $\pi^{*}E=\bigoplus_{g\in\ker\pi}\tau_{g}^{*}L$ is a direct
sum of line bundles of one slope, hence $\mu$-semistable, and semistability
descends along the finite \'etale $\pi$; so $E$ is semistable. If $F\subset E$
is a saturated subsheaf of the same slope, then $\pi^{*}F\subset\pi^{*}E$ is
a saturated subsheaf of the same slope of a direct sum of pairwise
non-isomorphic line bundles, the $\tau_{g}^{*}L$ being distinct because the
block is simple, hence a partial sum $\bigoplus_{g\in S}\tau_{g}^{*}L$; it is
invariant under the translations by $\ker\pi$, which permute the summands
transitively, so $S$ is empty or everything and $F=0$ or $F=E$. Finally, a
nonzero morphism $f$ between $\mu$-stable bundles of the same slope is
injective and generically surjective, so the two bundles have the same rank;
for blocks of the same slope $\nu$ they then have the same $c_{1}$, since
$c_{1}(E)=\nu\,\rk(E)\,\Theta$, and $\det f$ is a nonzero section of a
numerically trivial line bundle, hence nowhere zero. So $f$ is an
isomorphism.
\end{proof}

\begin{theorem}[No two-term complex of simple blocks satisfies the criterion]
\label{thm:noblocks}
Let $(X,\Theta)$ be a polarised abelian variety of dimension $n\ge3$, let $A$
and $B$ be direct sums of simple isogeny blocks, let $\varphi\colon A\to B$
be any morphism and let $F$ be the cone of $\varphi$ or a shift of it. If
$F\ne0$ and $\ch(F)=au_{\Theta}+bv_{\Theta}$ with $(a,b)\ne(0,0)$, then $F$
does not satisfy the weakened criterion. In particular, no object with a
resolution $0\to E_{1}\to E_{0}\to F\to0$ by sums of simple isogeny blocks
satisfies it, whatever the rank.
\end{theorem}

\begin{proof}
Cancel, by Gaussian elimination in the two-term complex, every pair of
isomorphic blocks of $A$ and $B$ joined by an entry of $\varphi$ that is an
isomorphism; this changes neither the cone nor the block structure of the
remaining terms. By \Cref{lem:blocks}(iii), every remaining entry between
blocks of the same slope is zero. If $A$ or $B$ is empty after cancelling,
$F$ is a shifted sum of blocks, on each of which $\operatorname{ev}(x_{\pi})$
is a nonzero scalar by (i). Otherwise the rest is the proof of
\Cref{thm:nosplit} with \Cref{lem:blocks} in place of the vanishing of
$H^{2}(M^{\otimes m})$. For a block $E'$ of $B$ with inclusion $\iota$ and
$q'=q\circ\iota\colon E'\to F$, the class $z\cdot\iota\in\Hom(E',B[2])$ has
component $z\cdot\id_{E'}\ne0$ in $\Ext^{2}(E',E')=H^{2}(\cO_{X})\cdot\id_{E'}$
by (ii). It lies in the image of $\varphi_{*}$ only if
$\sum_{i}\varphi_{ji}\psi_{i}=z\cdot\id_{E'}$ for some
$\psi_{i}\in\Ext^{2}(E',A_{i})$, and $\psi_{i}\ne0$ forces $A_{i}$ to have
the slope of $E'$ by (ii), for which $\varphi_{ji}=0$. So $z\cdot q'\ne0$,
contradicting \Cref{lem:annihilation} in the form given by (i).
\end{proof}

\begin{remark}[Semi-homogeneous bundles]\label{rem:blocks}
By Mukai \cite{Muk78}, the simple isogeny blocks are exactly the simple
semi-homogeneous bundles, a bundle $E$ being semi-homogeneous if for every
$x$ there is $P\in\Pic^{0}(X)$ with $\tau_{x}^{*}E\cong E\otimes P$; and
every semi-homogeneous bundle on a variety with $\NS(X)_{\QQ}=\QQ\Theta$ is
an iterated extension of such blocks and their twists by $\Pic^{0}$. We do
not use this. In these terms, \Cref{thm:noblocks} says that the bundles in a
Hilbert--Burch resolution of a candidate cannot all have scalar Atiyah
class. The obstruction is always the same: the compensated Poisson classes
act on a block by the scalar $\tfrac{\nu^{2}+d}{2}\,\pi\lrcorner\,\Theta^{2}$,
which would vanish only for $\nu^{2}=-d$, and cones of maps between blocks
inherit the scalar through the summand inclusions, because $\Ext^{2}$
between blocks of different slopes vanishes. A candidate must therefore be
built from a bundle whose Atiyah class is not a scalar, that is, one which
is not projectively flat.
\end{remark}

\begin{corollary}[Split resolutions and the remaining shapes]
\label{cor:disjointroutes}
Let $(X,\Theta)$ be a principally polarised abelian fourfold.
\begin{enumerate}[label=\textup{(\roman*)},nosep]
\item No secant object of rank one whose support has a resolution
  \eqref{eq:burchres} by sums of powers of $\Theta$ satisfies the weakened
  criterion, for any $r$, any twists and any $d$, whether or not its matrix
  degenerates in codimension two.
\item If a secant object $\cI_{Z}(b\Theta)$ satisfies the criterion, then
  $Z$ has no point at which two smooth branches of codimension two meet
  transversally, no two-term resolution of $\cI_{Z}$ splits into powers of
  $\Theta$ or into simple semi-homogeneous bundles, and the restriction
  $H^{2}(\cO_{X})\to H^{2}(\cO_{Z})$ is injective. If $Z$ is
  Cohen--Macaulay, $\cI_{Z}$ has projective dimension one. If $b=3$ and $Z$
  is smooth, then $d\le9$, so $d\in\{1,3,5,7\}$ when $N=(9+d)/2$ is an
  integer as in \Cref{setup:markman}, the two terms of a resolution do not
  both have classes in $\cL_{\Theta}$, and when $X$ is simple the
  injectivity is automatic.
\item At $b=3$, smooth supports and split resolutions are disjoint, for
  every rank and every twist of the resolution: a smooth support forces
  $d\le9$, and a resolution by bundles with classes in $\cL_{\Theta}$ forces
  $d\equiv15,23\pmod{24}$, hence $d\ge15$, so that, even without (i), no
  support could have both properties.
\end{enumerate}
\end{corollary}

\begin{proof}
Part (i) is \Cref{thm:nosplit} with $M=\cO_{X}(\Theta)$, the resolution
\eqref{eq:burchres} presenting $\cI_{Z}(b\Theta)$ as the cone of
$\varphi\otimes\cO(b\Theta)$. In (ii), the branches are excluded by
\Cref{cor:cmdichotomy}(i), projective dimension one comes from
\Cref{cor:cmdichotomy}(ii), the non-splitting from \Cref{thm:nosplit} and
\Cref{thm:noblocks}, the injectivity from the last assertion of
\Cref{lem:annihilation}, the bound on $d$ from
\Cref{cor:fourdiscriminants}, the statement on the two terms from
\Cref{cor:latticeburch}, and the automatic injectivity from
\Cref{lem:smoothinjective}. Part (iii) compares
\Cref{cor:fourdiscriminants} with \Cref{cor:latticeburch}: no $d\le9$ is
congruent to $15$ or $23$ modulo $24$.
\end{proof}
